\subsection{Librerías estándar útiles para ICPC}

Python posee una biblioteca estándar muy completa. Conocer estos módulos permite resolver muchos problemas sin implementar estructuras o algoritmos desde cero.

\vspace{0.3cm}

\subsubsection{collections}

Contiene estructuras de datos optimizadas.

\begin{lstlisting}[language=Python]
from collections import deque
from collections import defaultdict
from collections import Counter
from collections import OrderedDict
from collections import namedtuple
\end{lstlisting}

\begin{center}
\begin{tabular}{|l|p{6cm}|}
\hline
Clase & Uso \\ \hline
deque & Cola doble, BFS, Sliding Window \\ \hline
defaultdict & Dict con valor por defecto \\ \hline
Counter & Conteo automático de frecuencias \\ \hline
OrderedDict & Dict que conserva operaciones de orden \\ \hline
namedtuple & Tuplas con nombres para los campos \\ \hline
\end{tabular}
\end{center}

\vspace{0.3cm}

\paragraph{deque}

Permite insertar y eliminar elementos por ambos extremos en \bigO{1}

\begin{lstlisting}[language=Python]
from collections import deque
q = deque()
q.append(5)
q.appendleft(2)
q.pop()
q.popleft()
q.rotate(1)
q.rotate(-1)
\end{lstlisting}

Es la estructura estándar para implementar BFS.

\vspace{0.3cm}

\paragraph{defaultdict}

Ya estudiado en la sección anterior.

Se utiliza principalmente para:

\begin{itemize}
\item frecuencias;
\item listas de adyacencia;
\item agrupaciones.
\end{itemize}

\vspace{0.3cm}

\paragraph{Counter}

Muy útil para conteo de frecuencias.

\begin{lstlisting}[language=Python]
freq = Counter(arr)
freq.most_common(3)
\end{lstlisting}

\vspace{0.3cm}

\subsubsection{heapq}

Implementa un montículo mínimo (Min Heap).

\begin{lstlisting}[language=Python]
import heapq
\end{lstlisting}

Funciones principales.

\begin{center}
\begin{tabular}{|l|l|}
\hline
Función & Uso \\ \hline
heappush() & Insertar \\ \hline
heappop() & Extraer mínimo \\ \hline
heapify() & Construir heap \\ \hline
heapreplace() & Reemplazar mínimo \\ \hline
nlargest() & Mayores elementos \\ \hline
nsmallest() & Menores elementos \\ \hline
\end{tabular}
\end{center}

Ejemplo.

\begin{lstlisting}[language=Python]
heap=[]
heapq.heappush(heap,4)
heapq.heappush(heap,2)
heapq.heappush(heap,8)
heapq.heappop(heap)
\end{lstlisting}

Todas las operaciones de inserción y extracción tienen complejidad \bigO{log n}

Se utiliza principalmente en:

\begin{itemize}
\item Dijkstra;
\item Prim;
\item Huffman;
\item Top-K.
\end{itemize}

\vspace{0.3cm}

\subsubsection{bisect}

Permite realizar búsqueda binaria sobre listas ordenadas.

\begin{lstlisting}[language=Python]
from bisect import bisect_left
from bisect import bisect_right
from bisect import insort
\end{lstlisting}

\begin{lstlisting}[language=Python]
bisect_left(arr,x)
bisect_right(arr,x)
insort(arr,x)
\end{lstlisting}

Todas las búsquedas tienen complejidad \bigO{log n}

\vspace{0.3cm}

\subsubsection{functools}

Funciones de orden superior.

\begin{lstlisting}[language=Python]
from functools import lru_cache
from functools import cache
from functools import cmp_to_key
from functools import reduce
\end{lstlisting}

\paragraph{lru\_cache}

Muy utilizado para memoización.

\begin{lstlisting}[language=Python]
@lru_cache(None)
def dp(i,j):
    ...
\end{lstlisting}

\paragraph{cmp\_to\_key}

Permite utilizar comparadores personalizados.

\begin{lstlisting}[language=Python]
sorted(arr,key=cmp_to_key(cmp))
\end{lstlisting}

\vspace{0.3cm}

\subsubsection{itertools}

Herramientas para trabajar con iteradores.

\begin{lstlisting}[language=Python]
from itertools import permutations
from itertools import combinations
from itertools import combinations_with_replacement
from itertools import product
from itertools import accumulate
from itertools import cycle
from itertools import repeat
from itertools import pairwise
\end{lstlisting}

\begin{center}
\begin{tabular}{|l|l|}
\hline
Función & Uso \\ \hline
permutations & Permutaciones \\ \hline
combinations & Combinaciones \\ \hline
product & Producto cartesiano \\ \hline
accumulate & Prefijos acumulados \\ \hline
cycle & Ciclo infinito \\ \hline
repeat & Repetición \\ \hline
pairwise & Pares consecutivos \\ \hline
\end{tabular}
\end{center}

\vspace{0.3cm}

Ejemplos.

\begin{lstlisting}[language=Python]
permutations(arr)
combinations(arr,2)
product(a,b)
accumulate(arr)
\end{lstlisting}

\vspace{0.3cm}

\subsubsection{operator}

Contiene operadores como funciones.

\begin{lstlisting}[language=Python]
from operator import itemgetter
from operator import attrgetter
\end{lstlisting}

Ejemplo.

\begin{lstlisting}[language=Python]
sorted(arr,key=itemgetter(1))
\end{lstlisting}

\vspace{0.3cm}

\subsubsection{decimal}

Permite realizar aritmética decimal exacta.

\begin{lstlisting}[language=Python]
from decimal import Decimal
\end{lstlisting}

Normalmente no se utiliza en ICPC debido a su menor rendimiento.

\vspace{0.3cm}

\subsubsection{fractions}

Permite trabajar con fracciones exactas.

\begin{lstlisting}[language=Python]
from fractions import Fraction
\end{lstlisting}

Muy útil cuando se desea evitar errores de punto flotante.

\vspace{0.3cm}

\subsubsection{array}

Implementa arreglos homogéneos.

\begin{lstlisting}[language=Python]
from array import array
\end{lstlisting}

Consumen menos memoria que una lista convencional.

\vspace{0.3cm}

\subsubsection{queue}

\begin{lstlisting}[language=Python]
from queue import PriorityQueue
\end{lstlisting}

Aunque implementa colas de prioridad, en programación competitiva normalmente se prefiere utilizar

\begin{lstlisting}[language=Python]
heapq
\end{lstlisting}

por ser considerablemente más rápido.

\vspace{0.3cm}

\subsubsection{sys}

El módulo más utilizado para entrada y salida rápida.

\begin{lstlisting}[language=Python]
import sys

input = sys.stdin.readline

sys.setrecursionlimit(200000)
\end{lstlisting}

\vspace{0.3cm}

\subsubsection{Aplicaciones frecuentes en ICPC}

\begin{center}
\begin{tabular}{|l|l|}
\hline
Problema & Librería recomendada \\ \hline
BFS & deque \\ \hline
Dijkstra & heapq \\ \hline
Memoización & lru\_cache \\ \hline
Conteo & Counter \\ \hline
Frecuencias & defaultdict \\ \hline
Binary Search & bisect \\ \hline
Permutaciones & itertools \\ \hline
Combinaciones & itertools \\ \hline
Ordenamientos personalizados & cmp\_to\_key \\ \hline
Prefijos & accumulate \\ \hline
\end{tabular}
\end{center}

\vspace{0.3cm}

\subsubsection*{Resumen}

\begin{center}
\begin{tabular}{|l|l|}
\hline
Módulo & Uso principal \\ \hline
collections & Estructuras de datos \\ \hline
heapq & Heap mínimo \\ \hline
bisect & Búsqueda binaria \\ \hline
functools & Memoización y comparadores \\ \hline
itertools & Iteradores avanzados \\ \hline
operator & Funciones auxiliares para ordenar \\ \hline
fractions & Fracciones exactas \\ \hline
decimal & Decimales exactos \\ \hline
array & Arreglos homogéneos \\ \hline
queue & Colas sincronizadas \\ \hline
sys & Entrada rápida y recursión \\ \hline
\end{tabular}
\end{center}